\section{A monoid of diagrams}\label{sec:monoid}

The forcing construction will compare a diagram with its translated
subdiagrams.  We first construct the monoid that indexes these
translations.  Its two extension operations serve different purposes:
one realizes finitely many prescribed overlaps, and the other inserts
many new positions with a prescribed set of old divisors.  A further
property of the construction will control symmetries along countable
systems of subdiagrams.

All constructions in this section take place in a universe satisfying
$\ZFC$.  No inaccessible cardinal is used.  Fix an infinite cardinal
$\mu\geq\aleph_1$ such that $\mu^{\aleph_0}=\mu$.

\subsection{The required algebra and the construction theorem}

Let $M$ be a monoid with identity $1$. Write $x\sqsubseteq y$ if $xz=y$ for
some $z$.  We use only left cancellation.  A monoid is \emph{conical}
if $xy=1$ implies $x=y=1$.  Left cancellation and conicality make
$\sqsubseteq$ a partial order and make the right factor in a comparison
unique.  Write $x\sqsubset y$ for $x\sqsubseteq y$ and $x\ne y$,
and $\downarrow a=\{x:x\sqsubseteq a\}$.
A \emph{conditional join} is an element $x\vee y$ satisfying
\[
 xM\cap yM=(x\vee y)M
\]
whenever the intersection is nonempty.  Thus a join is required exactly
when the two elements have a common right multiple.
An \emph{ideal} in $\downarrow a$ will mean a nonempty downward closed
subset closed under finite nonempty joins.  These joins exist because
$a$ is an upper bound.

An inclusion $B\subseteq M$ is \emph{conservative} if it preserves and
reflects comparisons, preserves their right factors, preserves existing
joins, and preserves the absence of a common right multiple.  In
particular, if $x,y\in B$ and $xz=y$ in $M$, then $z\in B$ and the
same equation holds in $B$.  Mere injectivity is not this property.

We shall also keep real-valued height and cost functions $\rho,\epsilon$
on $M$, writing $\epsilon_x=\epsilon(x)$, and satisfying
\begin{equation}\label{eq:mono-margin}
 \begin{gathered}
 0<\rho(x)<1,\qquad 0<\epsilon_x\leq1,\\
 \rho(y)-\rho(x)>
 \sup\{\epsilon_d:d\sqsubseteq y,\ d\not\sqsubseteq x\}
 \qquad(x\sqsubset y).
 \end{gathered}
\end{equation}
An empty supremum of costs is understood to be zero.

\begin{lemma}\label{lem:mono-countable-ideal}
If a monoid has conditional joins and satisfies
Equation~\eqref{eq:mono-margin}, every ideal in $\downarrow a$ has a
nondecreasing cofinal sequence of length $\omega$, with repetitions
allowed.  The same conclusion holds in any $\ZFC$ extension in which
the monoid, order, heights and costs retain their meanings, for ideals
belonging to that extension.
\end{lemma}
\begin{proof}
For an ideal $U$, choose elements whose heights approach
$\sup\rho[U]$, and replace their successive initial segments by their
finite joins.  This gives a nondecreasing sequence $(n_i)$ in $U$
with heights approaching the same supremum.  If $d\in U$ divides
none of the $n_i$, then
\[
 \rho(n_i\vee d)-\rho(n_i)>\epsilon_d
\]
for every $i$, contradicting convergence to the supremum.  Prepending
$1$ makes $n_0=1$.  The proof applies in the stated extension as well;
its choices are made in that $\ZFC$ universe.
\end{proof}

Here is the finite overlap condition used below.  Take finitely many
copies $\{i\}\times M$.  For each pair of distinct copies prescribe
either no overlap or elements $a_{ij},a_{ji}\in M$ and the
identifications
\[
 (i,a_{ij}x)\sim(j,a_{ji}x)\qquad(x\in M).
\]
We call the prescription \emph{valid} if these pair identifications,
together with equality, already form an equivalence relation and do
not identify distinct elements of any one copy.  In particular,
transitive closure is not allowed to create extra overlaps.

\begin{theorem}\label{thm:monoid}
There is a monoid $T$ of cardinality $\mu$, with a distinguished
$e\ne1$, having the following properties.
\begin{enumerate}
\item $T$ is left cancellative and conical, has conditional joins,
and has heights and costs satisfying Equation~\eqref{eq:mono-margin}.
\item Every valid finite overlap prescription on copies of $T$ is
realized by left translations $x\mapsto l_i x$ in $T$, with exactly
the prescribed identifications.
\item For $a,b\in T$ with $b\ne1$, and every ideal
$U\subseteq\downarrow a$, there are $\mu$ distinct elements $m$ with
\begin{equation}\label{eq:mono-probes}
 \begin{gathered}
 mb=ae,\qquad m\not\sqsubseteq a,\qquad
 \downarrow m\cap\downarrow a=U,\\
 n\vee m=ae\quad(n\sqsubseteq a,\ n\notin U).
 \end{gathered}
\end{equation}
The cones of distinct elements in this family intersect exactly in
$aeT$.  Moreover $|T\setminus eT|=\mu$.
\item For every countable subset of $T$ there is a countable
submonoid $T_0$ containing it and a continuous conservative tower
starting at $T_0$ and ending at $T$.  Each successor step has one of
the following properties:
\begin{enumerate}
\item no new element divides an old element;
\item every new element $w$ with an old right multiple has a least
old right multiple $g(w)$, and
\begin{equation}\label{eq:mono-gate-old}
 d\text{ old},\quad g(w)\sqsubseteq wd
       \quad\Longrightarrow\quad wd\text{ old}.
\end{equation}
\end{enumerate}
Here ``old'' means belonging to the base of that particular successor
step, and the comparisons may be tested in the final monoid.
\end{enumerate}
\end{theorem}

For diagrams indexed by $T$, the finite-merger clause supplies a common
larger diagram for finitely many copies with compatible overlaps.  The
probe clause supplies $\mu$ copies with one prescribed intersection
pattern: for $n\sqsubseteq a$, the cone $nT$ contains $mT$ exactly when
$n\in U$, and otherwise meets it in $aeT$.  The margin makes every such
ideal countably cofinal.  The countable-base clause serves a different
purpose: it will allow suitably extendible positions below a countable
path to be moved into a countable submonoid.  We formulate that descent
precisely at the end of the section and use it to bound the residual
symmetries of diagrams.

We prove the theorem through two extension lemmas.  They preserve more
than the order: keeping right factors and nonjoins fixed is what permits
the operations to be rearranged over a countable submonoid.

\subsection{The extension that inserts a probe}

Suppose $M$ has all the algebraic and margin properties above.
Fix $a,b,e\in M$, with $b,e\ne1$, and a nondecreasing path
\begin{equation}\label{eq:mono-path}
 1=n_0\sqsubseteq n_1\sqsubseteq\cdots\sqsubseteq a,
 \qquad n_id_i=n_{i+1},\qquad n_ic_i=a.
\end{equation}
Left cancellation gives $d_ic_{i+1}=c_i$.
Form the monoid presentation over $M$ with new generators $p_i$ and
relations
\begin{equation}\label{eq:mono-probe-relations}
 p_i=d_ip_{i+1},\qquad p_ib=c_ie\qquad(i<\omega).
\end{equation}
Put $m=p_0$.  Since $n_0=1$ and $n_id_i=n_{i+1}$, the first
relations give $n_ip_i=m$ for every $i$.  Thus every point of the path
divides the new probe.  The second relations make both factorizations
in Figure~\ref{fig:probe-divisibility} end at $ae$:
\[
 n_ip_ib=n_ic_ie=ae,\qquad mb=ae.
\]
These identities explain the presentation, but do not yet establish
its exact divisor profile.  We must prove that no old divisor of $a$
outside the path's downward closure is inserted below $m$, that $m$
does not divide $a$, and that every omitted divisor meets $m$ first
at $ae$.  The next lemma proves these
claims even for new divisors of $a$ introduced by the same extension.

\begin{figure}[htbp]
\centering
\begin{tikzpicture}[font=\small,>=Stealth]
  \node (pathpoint) at (0,0) {$n_i$};
  \node (oldpoint) at (3,1) {$a$};
  \node (probe) at (3,-1) {$m=p_0$};
  \node (commonmultiple) at (6,0) {$ae$};
  \draw[->] (pathpoint) -- node[above left] {$c_i$} (oldpoint);
  \draw[->] (pathpoint) -- node[below left] {$p_i$} (probe);
  \draw[->] (oldpoint) -- node[above right] {$e$} (commonmultiple);
  \draw[->] (probe) -- node[below right] {$b$} (commonmultiple);
\end{tikzpicture}
\caption{The relations defining the probe give two factorizations of
the prescribed common multiple $ae$.  Each arrow means right
multiplication by its label, not a covering relation.  This is not a
Hasse diagram: path points may coincide with $a$, and $a$ may divide
$m$.  The exact-profile assertion is that a divisor of $a$ omitted
from the path's downward closure has join $ae$ with $m$.}
\label{fig:probe-divisibility}
\end{figure}


\begin{lemma}\label{lem:mono-probe-algebra}
The presented extension $M'$ is left cancellative and conical, has
conditional joins, and conservatively contains $M$.  Every new element
with an old right multiple has a least old right multiple, and
Equation~\eqref{eq:mono-gate-old} holds at this extension.
For $m=p_0$,
\begin{equation}\label{eq:mono-probe-profile}
 \begin{gathered}
 mb=ae,\qquad m\not\sqsubseteq a,\qquad
 \downarrow m\cap\downarrow a=\bigcup_i\downarrow n_i,\\
 x\vee m=ae\qquad
 (x\sqsubseteq a,\ x\not\sqsubseteq m),
 \end{gathered}
\end{equation}
where the downsets and joins in this display are computed in $M'$.
\end{lemma}

\begin{proof}
\smallskip\noindent\textit{Normal forms.}
Let $S$ be the direct limit of the sets $M$ under the maps
$x\mapsto xd_i$.  Write $[x,i]$ for the image of $x$ at stage $i$.
Equality in this direct limit means equality at a finite common stage.
The transition maps need not be injective: no right cancellation has
been assumed.  Nevertheless the left $M$-action on $S$ is injective
for each multiplier.  Indeed, equality after multiplying two letters
by an old $u$ becomes $uX=uY$ at a common stage, where left
cancellation gives $X=Y$.

There are well-defined equivariant maps
\[
 k([x,i])=xc_i,\qquad h(s)=k(s)e.
\]
The identity $d_ic_{i+1}=c_i$ proves well-definedness.  Before imposing
the second relation in Equation~\eqref{eq:mono-probe-relations}, the
normal forms are strings
\[
 s_1\cdots s_j y\qquad(s_i\in S,\ y\in M),
\]
including $j=0$.  The letter $[x,i]$ represents $xp_i$.
Old factors between new generators are absorbed into the following
letter, and the transition relations are precisely equality in $S$.
Multiplication concatenates strings, using the left action to absorb
the terminal old factor of the first string.

The remaining reductions are
\begin{equation}\label{eq:mono-reduction}
 s(bt)\longmapsto h(s)t,
\end{equation}
where $t$ is a letter of $S$ or a terminal old element.  The quotient
$t$ is unique because the left $b$-action is injective.  Each reduction
removes one new letter.  The only overlapping reductions occur in
$s\,(bt)\,(bu)$, with the final factor possibly terminal.  Reducing
in either order gives the same result by
\[
 h(bt)=b h(t),\qquad h(h(s)t)=h(s)h(t).
\]
Disjoint reductions commute by associativity of the action.  By Newman's
lemma \cite{Newman1942} (see \cite[Lemma~5.13]{DershowitzPlaisted2001}), these
local diamonds and termination give unique irreducibles: besides the
old elements, they are strings whose first letter is arbitrary and
whose later letters and terminal factor are not left-divisible by $b$.

Old left multipliers act injectively on these forms.  Left multiplication
by a letter $s$ either prepends $s$, when the input is not a $b$-multiple,
or sends $bZ$ to $h(s)Z$.  Both cases are injective.  Their ranges are
disjoint, because
\begin{equation}\label{eq:mono-head-separation}
 s\notin h(s)S.
\end{equation}
To see this, $s\in xS$ implies $x\sqsubseteq k(s)$, whereas
$k(s)e\not\sqsubseteq k(s)$ by conicality and left cancellation.
Thus every left multiplication is injective.  A new word can become
old only by consuming its first letter, at which point it is a multiple
of $h(s)\ne1$.  It cannot equal $1$, proving conicality.

\smallskip\noindent\textit{Old cones and the first-letter blocks.}
For $s\in S$, let $B_s$ consist of all irreducibles beginning with $s$.
For old $x$ the normal forms give
\begin{equation}\label{eq:mono-old-cone}
 xM'=xM\ \cup\!\bigcup_{s\in xS}B_s.
\end{equation}
At a common stage in the direct limit,
$xS\cap yS$ is $(x\vee y)S$ if the old join exists, and is empty
if the old cones are disjoint.  Thus old comparisons, their right
factors, joins and nonjoins are unchanged.  Notice that this conclusion
has not used a join theorem for arbitrary new words.

\smallskip\noindent\textit{Exits and gates.}
Fix an irreducible $w=sY$.  Right multiplication either leaves its
first letter intact or eventually consumes that letter.  We prove the
following precise alternative: the set $wM'\setminus B_s$ is empty,
or it is one cone $g(w)M'$ whose root $g(w)$ is old.  Thus the gate
will describe every exit from the fixed-head block, not only the old
multiples of $w$.

The first letter is consumed exactly when the suffix becomes a
$b$-multiple, so first calculate $YM'\cap bM'$.  If $Y$ is old,
Equation~\eqref{eq:mono-old-cone} gives the old join $Y\vee b$,
or an empty intersection.  If $Y$ is new with head $t$,
irreducibility gives $t\notin bS$.  No member of $B_t$ is then a
$b$-multiple.  By induction on the number of letters in the suffix,
all possible exits of $Y$ form the old-rooted cone $g(Y)M'$, if
there are any.  Its intersection with $bM'$ is calculated by the
old join $g(Y)\vee b$.

In either case the suffix intersection is empty or is $zM'$ for
an old $z=bu$ with old $u$.  The factor $u$ is old by conservation
of old comparisons.  In the nonempty case, multiplication by $s$
and the reduction $sbu=h(s)u$ give
\begin{equation}\label{eq:mono-gate-formula}
 wM'\setminus B_s=g(w)M',\qquad g(w)=h(s)u.
\end{equation}
This cone cannot return to $B_s$: every new head in it belongs to
$h(s)S$, whereas $s\notin h(s)S$ by
Equation~\eqref{eq:mono-head-separation}.  This proves the exit
alternative.  Every old multiple of $w$ lies outside $B_s$, so
$g(w)$ is in particular its least old multiple.

We record also the right factor $c(w)$ with $wc(w)=g(w)$.
If $Y$ is old, it is the old quotient of $z$ by $Y$.
If $Y$ is new and $Yc(Y)=g(Y)$, it is
\begin{equation}\label{eq:mono-gate-quotient}
 c(w)=c(Y)v,
 \qquad g(Y)v=z.
\end{equation}
All quotients in this recursion are old except the already constructed
$c(Y)$.  This explicit calculation will ensure conservation of right
factors when the same presentation is simulated over a larger base.

Append an old $d$ to $w$.  If the product leaves $B_s$, its first
letter has been consumed; that reduction can occur only after all
intervening letters have been consumed.  The result is then entirely
old.  Since the gate cone is disjoint from $B_s$, the hypothesis
$g(w)\sqsubseteq wd$ implies just such an exit.  This proves
Equation~\eqref{eq:mono-gate-old}.

\smallskip\noindent\textit{All conditional joins.}
For two forms with the same head, cancel that head and intersect their
suffix cones, using induction on the total number of new letters.
For old $x$ and new $w\in B_s$, either $s\in xS$, in which case
$x\sqsubseteq w$, or every common multiple must leave $B_s$.
The latter intersection is computed through $g(w)$ and an old join.
For $w\in B_s$ and $v\in B_t$ with $s\ne t$, an intersection
inside $B_s$ requires $g(v)$ to divide a member of $B_s$.  By
Equation~\eqref{eq:mono-old-cone} it then divides all of $B_s$, so
$v\sqsubseteq w$.  The symmetric case is the same.  Otherwise a
common multiple must leave both blocks, and the intersection is
computed by the old join of their gates.  Missing gates or missing
old joins give an empty intersection.  This proves conditional joins
in every case.

\smallskip\noindent\textit{The prescribed divisor profile.}
For any letter $s$, the old divisors of every element of $B_s$ form
\begin{equation}\label{eq:mono-letter-ideal}
 P_s=\{x\in M:s\in xS\}\subseteq\downarrow k(s).
\end{equation}
This is an ideal: common-stage representatives and old joins prove
join closure.  For the head $[1,0]$ of $m=p_0$, it is precisely
$\bigcup_i\downarrow_M n_i$, and the gate is $ae$.  The defining
relation gives $mb=ae$.  Since $ae\not\sqsubseteq a$, no member
of $B_{[1,0]}$ divides $a$; in particular $m\not\sqsubseteq a$.

Now let $x$ be a new divisor of $a$.  It is outside that block and
has an old gate $g(x)\sqsubseteq a$.  If $x\sqsubseteq m$, the
comparison leaves its block, so $g(x)\sqsubseteq m$, and the old
profile puts $g(x)$ below some $n_i$.  Conversely $x\sqsubseteq n_i$
implies $x\sqsubseteq m$.  If $x\not\sqsubseteq m$, its gate also
fails to divide $m$.  The possibility $m\sqsubseteq x$ is excluded
by $x\sqsubseteq a$.  The gate calculation consequently gives
$x\vee m=g(x)\vee m=ae$.  The corresponding assertions for old
$x$ follow directly from its block membership.  This proves
Equation~\eqref{eq:mono-probe-profile}, including new divisors of $a$.
\end{proof}

\subsection{Extending the real heights}

The gate description also controls the costs of all newly introduced
divisors.  This is the reason for keeping a cost as well as a height:
pointwise strict inequalities alone would not control suprema at a
limit stage.

\begin{lemma}\label{lem:mono-probe-heights}
The height and cost functions on $M$ extend to the probe extension
$M'$ so that Equation~\eqref{eq:mono-margin} holds.  For any old
comparison, its missed-cost supremum is unchanged.
\end{lemma}
\begin{proof}
Use $P_s$ from Equation~\eqref{eq:mono-letter-ideal}, and put
$A_s=\sup\rho[P_s]$.  We place the whole block $B_s$ in a short
height interval above its old divisors and below $\rho(h(s))$.
Here $h(s)$ is the gate of the bare letter $s$; the gate of a longer
word in $B_s$, when present, is a multiple of $h(s)$ and need not
equal it.  The interval must leave enough room for the costs of both
the old divisors and the new words.  More precisely, there are
positive $u_s,\delta_s$ with
\begin{equation}\label{eq:mono-height-interval}
 \begin{split}
 u_s-A_s&>\delta_s,\\
 \rho(h(s))-u_s-\delta_s
  &>\sup\{\epsilon_d:d\sqsubseteq h(s),\ d\in M\setminus P_s\}.
 \end{split}
\end{equation}
Indeed, put $p=\rho(k(s))-A_s\geq0$ and
$q=\rho(h(s))-\rho(k(s))>0$.  For $d\sqsubseteq k(s)$ missing
$P_s$, every $x\in P_s$ misses $d$, and the old margin yields
$\epsilon_d<\rho(k(s))-\rho(x)$.  The supremum of these costs is
at most $p$.  The supremum for $d\sqsubseteq h(s)$ not dividing
$k(s)$ is strictly below $q$.  Thus the full supremum on the right
of Equation~\eqref{eq:mono-height-interval} is strictly less than
$p+q=\rho(h(s))-A_s$, also when $p=0$.  This leaves the required
positive room for $u_s$ and $\delta_s$.

Define recursively on irreducibles
\begin{equation}\label{eq:mono-probe-height-definition}
 \begin{gathered}
 \rho(sY)=u_s+\delta_s\rho(Y),\qquad
 0<\epsilon_{sY}\leq\delta_s\epsilon_Y,\\
 \epsilon_{sY}\leq\epsilon_{g(sY)}
       \quad\text{whenever the gate exists}.
 \end{gathered}
\end{equation}
These choices give heights in $(0,1)$ and costs in $(0,1]$.
We verify all comparisons, including the suprema.

For two elements in one block $B_s$, a divisor of the upper element
from outside that block also divides the lower one.  An old divisor
belongs to $P_s$; a new divisor with a different head must enter by
its old gate, which then divides every member of $B_s$.
Inside the block, cancellation removes $s$, and the margin follows
by induction from the factor $\delta_s$ in both height differences
and cost bounds.

Suppose an old $x$ lies strictly below $sY$.  Then $x\in P_s$.
A missed divisor outside $B_s$ is old in $P_s$, or has an old gate
there of at least its cost.  For such an old divisor $d$ missing
$x$, the join $x\vee d$ lies in $P_s$, and hence
\[
 \epsilon_d<\rho(x\vee d)-\rho(x)\leq A_s-\rho(x).
\]
Their supremum is at most $A_s-\rho(x)$.  Divisors inside $B_s$
have cost at most $\delta_s$.  Equation~\eqref{eq:mono-height-interval}
makes $\rho(sY)-\rho(x)$ strictly larger than both bounds.

Suppose $sY$ lies strictly below an old $y$.  Then
$h(s)\sqsubseteq y$.  Every divisor of $y$ missing $sY$ is old
outside $P_s$, or has a gate there with at least its cost: a gate
in $P_s$ would put the divisor below all of $B_s$.
Among these old costs, those for divisors of $h(s)$ are bounded by
Equation~\eqref{eq:mono-height-interval}.  If $y\ne h(s)$, the
remaining costs have supremum strictly below
$\rho(y)-\rho(h(s))$ by the old margin.  The total gap
$\rho(y)-\rho(sY)$ is the sum of two positive gaps and exceeds both
cost suprema.  This proves this comparison case.

Distinct comparable new blocks have an intervening old gate, so the
preceding cases apply transitively.  To justify this use, if
$x<z<y$, any divisor of $y$ missing $x$ either divides $z$ and
misses $x$, or misses $z$; the sum of the two positive height gaps
strictly dominates both relevant suprema.

Finally, for old $x<y$, any new divisor $d$ of $y$ missing $x$
has an old gate $g(d)\sqsubseteq y$ missing $x$, with
$\epsilon_d\leq\epsilon_{g(d)}$.  No new cost increases the old
supremum, while every old divisor remains present.  The supremum is
therefore exactly unchanged.  This also proves persistence of the
margins at unions of a conservative tower of these extensions.
\end{proof}

\subsection{The extension that merges finitely many copies}

\begin{lemma}\label{lem:mono-merger}
Let $M$ satisfy the algebraic and margin hypotheses, and let a finite
nonempty overlap prescription on copies of $M$ be valid.  There is a
conservative extension realizing it by left translations.  The
extension is left cancellative and conical, has conditional joins,
and admits extending heights and costs satisfying
Equation~\eqref{eq:mono-margin}.  No new element divides an old one.
\end{lemma}
\begin{proof}
Let $D$ be the prescribed quotient of the copies, writing $t_i x$
for the image of $x$ in copy $i$.  This quotient carries one layer
of the required overlaps; strings of its elements will complete that
layer to a monoid.  It is a right $M$-set and each
map $x\mapsto t_i x$ is injective.  On $D$ use the order
$z\sqsubseteq w$ when $zq=w$ for some $q\in M$; left cancellation
and conicality inside a representing copy make this a partial order.
Its right cones have conditional joins.  Within one copy this is the
old join.  For elements represented in different copies, every common
multiple must lie in their specified overlap, if there is one.
In each copy join the element with the start of that overlap; transfer
the two resulting elements into the overlap and join there.  A missing
join at any step means that the original cones are disjoint.  This
calculates exactly their intersection, because every common multiple
must pass through both of those preliminary joins.

Adjoin generators $l_i$ with the specified relations
$l_i a_{ij}=l_j a_{ji}$.  The normal forms are
\begin{equation}\label{eq:mono-merger-forms}
 xz_1\cdots z_j\qquad(x\in M,\ z_i\in D),
\end{equation}
including $j=0$.  Each relation acts inside a new-generator/old-factor
pair, exactly as in $D$; it changes neither the number of new letters
nor the preceding old coefficient.  This proves uniqueness of the
forms.  Multiplication acts by the leading old coefficient of the
second word on the last $D$-letter of the first, and then appends the
remaining letters.  For a fixed $z=t_i a$, the map $y\mapsto zy$
is injective by left cancellation in copy $i$.  Together with old
left cancellation, this proves cancellation for arbitrary products of
forms.  Letter count proves conicality and that no new word divides
an old element.

Unlike the probe extension, this construction never consumes a new
letter.  Its cones are therefore read from the final letter rather
than from the first-letter blocks used above.  For a positive-length
word, right multiplication fixes the prefix before its final letter,
advances that letter within its $D$-cone, and possibly appends further
letters.  More explicitly, compare
$xz_1\cdots z_p$ and $yt_1\cdots t_q$ with $p,q>0$.
A common multiple requires $x=y$ and equality of the letters before
position $\min(p,q)$.  If $p=q$, their intersection is obtained by
joining $z_p$ and $t_p$ in $D$, or is empty if that join is absent.
If $p<q$, a common multiple exists exactly when $z_p$ divides $t_p$
in the right $M$-set $D$; then the first word divides the second and
the second is their join.  The case $q<p$ is symmetric.
An old $x$ and a positive-length word with head $y$ have a common
multiple exactly when $x\sqsubseteq y$, in which case $x$ divides
that word.  Two old words use their old join, because the heads of
common multiples are precisely their old common multiples.
These descriptions prove conditional joins and preserve old comparisons,
quotients and nonjoins.  The words $l_i x$ have exactly the
identifications in $D$, as required.

It remains to supply margins.  If the number of copies is $r$, set
\[
 \rho_D(z)=\frac1r\sum_{i<r}\rho_i(z),\qquad
 \rho_i(z)=
 \begin{cases}
 \rho(x),&z=t_i x,\\
 0,&z\text{ has no representation in copy }i.
 \end{cases}
\]
Each coordinate is nondecreasing along a comparison, and some
coordinate increases strictly in a proper comparison.  Choose positive
costs with
\[
 \epsilon^D_z\leq\frac1r
 \min\bigl(\{\rho(1)/2\}\cup
           \{\epsilon_x:z=t_i x\text{ for some }i\}\bigr).
\]
For $u<v$ in $D$, classify a divisor $d$ of $v$ missing $u$ by a
copy representing $d$.  That copy also represents $v$.  If it
represents $u$, its old margin bounds the supremum of these costs
strictly below its coordinate increase divided by $r$.  Otherwise
that increase is at least $\rho(1)$, while each such cost is at most
$\rho(1)/(2r)$.  There are only finitely many copies, so their strict
bounds still give the required strict supremum bound for
$\rho_D(v)-\rho_D(u)$.

Choose $\theta,\phi>0$ with $\theta+\phi<1$, and recursively set
\begin{equation}\label{eq:mono-merger-heights}
 \begin{split}
 \rho(wz)&=\rho(w)+(1-\rho(w))(\theta+\phi\rho_D(z)),\\
 0<\epsilon_{wz}&\leq
 (1-\rho(w))\min(\theta/2,\phi\epsilon^D_z).
 \end{split}
\end{equation}
For comparable new words of equal length only their last letters
vary; missed divisors of that length are handled by the scaled
$D$-margin, and all shorter divisors divide both words.  For a word
$w$ compared with $wz$, every missed divisor has the new length and
prefix $w$, so its cost is at most $(1-\rho(w))\theta/2$, strictly
less than the height increase.  An arbitrary comparison factors
through these comparisons, and, when the lower word is old, an old
comparison of leading coefficients.  Transitivity of the margin,
proved in Lemma~\ref{lem:mono-probe-heights}, finishes the check.
Old cost suprema are unchanged because there are no new old divisors.
\end{proof}

\subsection{Conservation under simulation}

We next establish the compatibility needed to reorder the construction.
This step also explains why overlap prescriptions remain valid when
operations are scheduled before they are executed.

\begin{lemma}\label{lem:mono-simulation}
Suppose $B\subseteq M$ is conservative and both monoids have the
algebraic properties above.  Either extension operation with parameters
in $B$ is valid over $M$, and the extension of $B$ conservatively
embeds into the corresponding extension of $M$ by the same generators.
A finite overlap prescription with parameters in $B$ is valid in $B$
if and only if it is valid in $M$.
\end{lemma}
\begin{proof}
For the last assertion, the only nontrivial transitivity check passes
through three copies.  In the middle copy, intersect the two prescribed
overlap starts.  If their old cones are disjoint, conservativity keeps
them disjoint.  Otherwise their old join and its right quotients are
unchanged.  Every later common multiple is a multiple of that join.
The two endpoint identifications at the join already agree in a valid
prescription over $B$, so right multiplication gives agreement at all
later common multiples.  This proves validity over $M$.  Conversely,
restrict a valid prescription to $B$.  Comparisons between its old
parameters have their quotients in $B$, so no identification or
transitivity witness between elements of $B$ requires leaving $B$.
The restriction has exactly the indicated pair relations and is valid.

For probes, equality of two letters in the smaller direct limit is
preserved and reflected because it is an equality at a common finite
stage.  If an old $b\in B$ divides a smaller letter $[x,i]$ in the
larger direct limit, there is a common stage at which
\[
 by=x d_i\cdots d_{j-1}.
\]
Both endpoints of this comparison belong to $B$, so its right factor
$y$ belongs to $B$.  Thus divisibility of a smaller letter by an old
smaller multiplier, including its quotient, is unchanged.  The same
irreducible words therefore embed.  The gate recursion uses only the
previous gates and old joins and quotients, so it has the same outcome,
including the absence of a gate.  The join calculations in
Lemma~\ref{lem:mono-probe-algebra} now preserve and reflect joins and
nonjoins.

For clarity, preservation of right quotients between new forms follows
by a separate division argument.  If both forms are old, use the base
hypothesis.  If the divisor is old and the dividend new, divide its
first direct-limit letter by the preceding common-stage calculation,
and retain its suffix.  If both forms have the same new head, cancel
that head and recurse on their suffixes.  In the remaining cases the
new divisor must leave its first-letter block.  Factor through its
gate, using the quotient in Equation~\eqref{eq:mono-gate-quotient},
and divide that old gate into the dividend by the first two cases.
Every factor obtained lies in the smaller extension.  Left
cancellation makes their product the unique required quotient.

For mergers, representation in an overlap and intersections of cones
in $D$ were calculated solely using old joins and their right
quotients.  These calculations are consequently unchanged.  The
normal forms in Equation~\eqref{eq:mono-merger-forms} then preserve
comparisons, their right factors, joins and nonjoins.  This proves
conservative simulation for both operations.
\end{proof}

\begin{lemma}\label{lem:mono-profile-persistence}
Suppose a probe satisfies Equation~\eqref{eq:mono-probe-profile} in
one stage of a tower of the two operations.  It continues to satisfy
that equation at every later stage and at the union, with the full
later downward closure of its original path.  Previously constructed
joins between probes are also preserved.
\end{lemma}
\begin{proof}
A merger introduces no new divisor of the old $a$.  At a probe
extension, a new $x\sqsubseteq a$ has an old gate
$g(x)\sqsubseteq a$.  If $x$ also divides the old probe, its gate
lies below that probe and hence below some original path element;
then $x$ does too.  If $x$ does not divide the old probe, its gate
does not either.  The old probe cannot divide $x$, because it does
not divide $a$.  The join calculation through the gate therefore
gives $ae$ in this case.  This proves both alternatives at a successor
step.  At a limit every tested element has already appeared, and all
comparisons and joins have been preserved.  Conservativity directly
preserves joins between existing probes.
\end{proof}

\subsection{The schedule and the countable-base rearrangement}

We now assemble the operations.  The height bound ensures that the
apparently arbitrary ideal requirements can be scheduled using only
countable paths.  The rearrangement proved afterward concerns the
same resulting monoid, not a separately constructed substitute.

\begin{proof}[Proof of Theorem~\ref{thm:monoid}]
\smallskip\noindent\textit{Saturating the geometric requirements.}
Start with the free monoid on $e$.  For $n<\omega$, for example, put
\[
 \rho(e^n)=1-2^{-n-1},\qquad \epsilon_{e^n}=2^{-n-3}.
\]
These satisfy Equation~\eqref{eq:mono-margin}.  Perform $\omega_1$
rounds.  At the start of a round schedule every available probe datum
$(a,b,(n_i))$ with the fixed $e$, and every available finite valid
merger.  For each probe datum schedule $\mu$ successive insertions.
Execute the resulting wellordered list with the two lemmas, and take
unions at limits.  Lemma~\ref{lem:mono-simulation} keeps every
scheduled merger valid until it is executed.

A monoid of size at most $\mu$ has at most $\mu^{\aleph_0}=\mu$
countable paths and at most $\mu$ finite prescriptions.  A probe
extension uses a quotient of $M\times\omega$ and finite strings;
a merger uses finitely many copies and finite strings.  Thus every
operation preserves size at most $\mu$, and each round has at most
$\mu$ steps.  The union of $\omega_1$ such rounds still has size
at most $\mu$.  The margin lemmas ensure that old values are never
changed and that old missed-cost suprema never increase, so the final
union $T$ has Equation~\eqref{eq:mono-margin} as well as all the
algebraic properties.

Given an ideal in $\downarrow_T a$, use
Lemma~\ref{lem:mono-countable-ideal} to obtain a cofinal path starting
at $1$.  Its elements, $a,b$, and all its transition and terminal
quotients occur before one round: the set is countable and the rounds
have cofinality $\omega_1$.  Alternatively, quotient conservativity
already puts every quotient into any stage containing its endpoints.
The corresponding requirement is scheduled.  Its final lower closure
is the given ideal, and Lemma~\ref{lem:mono-profile-persistence}
therefore gives Equation~\eqref{eq:mono-probes} for the $\mu$
insertions.  For any earlier and later insertion with the same datum, the
earlier probe is old but is not a divisor of $a$, so it is not an old
divisor of the new probe.  Joining through the new probe's gate gives
$ae$.  This join persists in the final union, proving independence.

Every final finite overlap prescription has its parameters in one
stage.  Its validity reflects to that stage by
Lemma~\ref{lem:mono-simulation}, and the schedule realizes it later.
The realization remains exact by conservation of comparisons and
quotients: equality $l_i x=l_j y$ in a later stage must still factor
through the preserved join of the two roots, with its prescribed
quotients.  Thus later right factors add exactly the stipulated
identifications and no others.  Finally use $a=1,b=e$ and $U=\{1\}$.
Its probes have no old divisor $e$, so there are $\mu$ elements of
$T\setminus eT$.  This proves both cardinality assertions.

\smallskip\noindent\textit{Extracting a countable initial base.}
It remains to prove the local construction property for this same $T$.
We first select countably many operations containing any prescribed
countable set of elements, and then perform all omitted operations over
their union.  The resulting tower must retain the gate rule: this is what
will let a later descent replace a position by one born at an earlier
stage.  Regard the
entire schedule as a continuous tower $(M_\alpha)_{\alpha\leq\lambda}$
of individual operations.  Record each parameter by a finite word in
previous generators.  Every element of the presented monoid has such
a word, and each equality has a finite derivation from the relations.
This does not assert that a whole countable path has finite support.

Choose a sufficiently large regular $\Theta$ and a countable elementary
submodel $N$ of $H_\Theta$ containing the tower, its presentations and
an enumeration of the prescribed countable subset of $T$.
Countable objects in $N$, in particular the parameter paths of
operations indexed in $N$, have all their entries in $N$.
First perform just the operations whose indices belong to $N$, in
their original order.  Let $S_\alpha$ denote this subconstruction
using selected indices below the cut $\alpha$.

At cuts $\alpha\in N$ we have
\begin{equation}\label{eq:mono-elementary-cut}
 S_\alpha=M_\alpha\cap N,
\end{equation}
with a conservative inclusion in $M_\alpha$.  Here is the induction
justifying the assertion.  Old comparisons, their quotients and joins
in an elementary intersection are reflected by elementarity; a
nonjoin remains a nonjoin in the larger monoid.  Before a selected
operation, every parameter is already present: its finite term code
belongs to $N$, and all entries of a countable path in $N$ belong to
$N$.  Conservative simulation makes that operation legal in the
subconstruction and embeds its result.  Conversely, every element of
$M_\alpha\cap N$ has by elementarity a finite term code in $N$;
all of its finitely many generator indices belong to $N$ and precede
$\alpha$.  It is therefore produced by the selected subconstruction.
This proves Equation~\eqref{eq:mono-elementary-cut} and the successor
step of the induction.  At a cut $\alpha\notin N$, fix the finitely
many words in a comparison or join test.  If they use selected operation
indices, let $\gamma$ be the greatest such index and put
$\beta=\gamma+1$; if there are no such indices, put $\beta=0$.
Then $\beta\in N$, $\beta<\alpha$, and all the words belong to
$S_\beta=M_\beta\cap N$.  The original inclusion
$M_\beta\subseteq M_\alpha$ is conservative.  Any right quotient or
conditional join required by the test therefore already lies in
$M_\beta$ and, by elementarity, in $N$, hence in $S_\beta$.
This also preserves nonjoins and proves conservation at $\alpha$.
Unions introduce no additional finite relation.  This completes the
induction.

The resulting selected submonoid $T_0$ is countable: the initial
monoid is countable, there are countably many selected steps, and each
operation over a countable base adds only countably many elements.
It contains the prescribed sequence by its finite term codes.
The set of selected indices need not be an initial segment; the
argument has used Equation~\eqref{eq:mono-elementary-cut} only at
cuts belonging to $N$.

\smallskip\noindent\textit{Reinserting the omitted operations.}
Append now the omitted steps in their original order.  Before an
omitted step $\alpha$, the current presentation is the union of the
original relations below $\alpha$ and all selected relations.  It has
an alternative realization: start with $M_\alpha$, and then perform
the selected steps at indices at least $\alpha$ in their order.
Indeed $S_\alpha\subseteq M_\alpha$ is conservative, and repeatedly
simulating the selected subconstruction makes all these operations
legal.  Each extends its whole base conservatively, so this alternative
realization conservatively contains $M_\alpha$.  It has precisely
the same generators and relations as the current reordered
presentation.  Consequently the original operation at $\alpha$ is
legal over the current base, by conservative simulation again.
Induction validates every appended step.  At limits the presentations
are their unions; at the end their union is exactly the original
presentation of $T$.

Each step of this reordered tower is one of the two operations.
The merger has no new old divisors, and the probe extension has its
least-old-multiple rule, including
Equation~\eqref{eq:mono-gate-old}.  Conservation allows those
comparisons to be tested in $T$.  This proves the last assertion and
completes the theorem.
\end{proof}

We isolate the consequence of the local construction property that will
be used for symmetries.  It makes explicit why a position below a
countable path can be moved to the countable base after finitely many
steps; it does not claim that every such position was in that base
already.

\begin{lemma}\label{lem:mono-descent}
Let $(n_i)$ be a nondecreasing sequence in $T$, and choose a
countable-base tower from Theorem~\ref{thm:monoid} whose initial
submonoid contains all $n_i$ and all right quotients between them.
Let $\mathcal A\subseteq\omega\times T$ be a set of pairs such that
each $(l,m)\in\mathcal A$ has $m\sqsubseteq n_N$ for some $N$.
Assume that for each $(l,m)\in\mathcal A$ and each $N$ there is
$(l',md)\in\mathcal A$ with $l'\geq l$, $n_ld=n_{l'}$, and
$n_N\sqsubseteq md$.  Then, starting from any pair in $\mathcal A$,
a finite succession of these replacements reaches a pair whose second
coordinate belongs to the initial countable submonoid.
\end{lemma}
\begin{proof}
If $m$ is not in that submonoid, take its first birth step in the
reordered tower.  It has the old multiple $n_N$, so this is a probe
step and its gate satisfies $g(m)\sqsubseteq n_N$.  Choose an
admissible replacement with $n_N\sqsubseteq md$.  The quotient
$d$ belongs to the initial submonoid and is therefore old at this
step.  Conservation and Equation~\eqref{eq:mono-gate-old} imply
that $md$ belongs to the strictly earlier base.  Repeat if necessary.
Each repetition strictly lowers the birth ordinal, so the process
terminates after finitely many repetitions.  The resulting position
lies in the countable initial submonoid.
\end{proof}
